
\documentclass{article}
\usepackage{colortbl}
\usepackage{makecell}
\usepackage{multirow}
\usepackage{supertabular}

\begin{document}

\newcounter{utterance}

\centering \large Interaction Transcript for game `bbh', experiment `formal\_fallacies', episode 226 with reflexion{-}llama{-}3.1{-}8b{-}instruct.
\vspace{24pt}

{ \footnotesize  \setcounter{utterance}{1}
\setlength{\tabcolsep}{0pt}
\begin{supertabular}{c@{$\;$}|p{.15\linewidth}@{}p{.15\linewidth}p{.15\linewidth}p{.15\linewidth}p{.15\linewidth}p{.15\linewidth}}
    \# & \multicolumn{2}{c}{Player} && \multicolumn{2}{c}{Game Master} \\
    \hline

    \theutterance \stepcounter{utterance}  
    & & & \multicolumn{4}{p{0.6\linewidth}}{
        \cellcolor[rgb]{0.9,0.9,0.9}{
            \makecell[{{p{\linewidth}}}]{
                \texttt{\tiny{[P1$\langle$GM]}}
                \texttt{You have attempted this same task before and did not succeed. These are the plans you wrote after those attempts.} \\
\texttt{Trial \#1: New plan: } \\
\\ 
\texttt{Upon reflection, I realize that I failed to identify the formal fallacies in the arguments. In particular, I did not recognize the invalidity of the third argument. To improve my strategy, I will focus on carefully analyzing the premises and the conclusion of each argument, looking for any logical flaws or inconsistencies. I will also pay closer attention to the use of quantifiers (e.g., "every," "some," "all") and the relationships between the different statements. Specifically, I will ask myself questions like: "Does the conclusion logically follow from the premises?" "Are there any implicit assumptions or unstated premises that could affect the validity of the argument?" "Are there any logical fallacies, such as affirming the consequent or denying the antecedent?" By asking these questions and carefully examining the arguments, I hope to improve my ability to distinguish deductively valid arguments from formal fallacies.} \\
\\ 
\texttt{Trial \#2: New plan: } \\
\\ 
\texttt{Upon reflection, I realize that I failed to identify the formal fallacies in the arguments in my previous attempts. In particular, I did not recognize the invalidity of the third argument. To improve my strategy, I will focus on carefully analyzing the premises and the conclusion of each argument, looking for any logical flaws or inconsistencies. I will also pay closer attention to the use of quantifiers (e.g., "every," "some," "all") and the relationships between the different statements. Specifically, I will ask myself questions like: "Does the conclusion logically follow from the premises?" "Are there any implicit assumptions or unstated premises that could affect the validity of the argument?" "Are there any logical fallacies, such as affirming the consequent or denying the antecedent?" By asking these questions and carefully examining the arguments, I hope to improve my ability to distinguish deductively valid arguments from formal fallacies.} \\
\\ 
\texttt{To apply this strategy more effectively, I will also pay closer attention to the structure of the arguments, including the use of conditionals, disjunctions, and negations. I will also try to identify any potential ambiguities or unclear statements in the premises or conclusion, and ask myself how these might affect the validity of the argument.} \\
\\ 
\texttt{In particular, I will focus on the third argument, which I previously failed to identify as invalid. Upon re{-}examining the argument, I realize that I made a mistake in my previous analysis. The correct analysis is as follows:} \\
\\ 
\texttt{The premises state that every infrequent user of Paul Mitchell shampoo is either a rare consumer of Nioxin shampoo or a loyal buyer of Caress soap, or both. However, the conclusion states that whoever is an infrequent user of Paul Mitchell shampoo is not a regular consumer of Lush soap. This conclusion does not logically follow from the premises, as there is no information that infrequent users of Paul Mitchell shampoo are not regular consumers of Lush soap. In fact, it is possible that some infrequent users of Paul Mitchell shampoo are also regular consumers of Lush soap. Therefore, the argument is invalid.} \\
\\ 
\texttt{By carefully analyzing the premises and conclusion of this argument, I hope to improve my ability to identify formal fallacies and distinguish deductively valid arguments from invalid ones.} \\
\\ 
\texttt{Trial \#3: New plan: } \\
\\ 
\texttt{Upon reflection, I realize that I failed to identify the formal fallacies in the arguments in my previous attempts. In particular, I did not recognize the invalidity of the third argument. To improve my strategy, I will focus on carefully analyzing the premises and the conclusion of each argument, looking for any logical flaws or inconsistencies. I will also pay closer attention to the use of quantifiers (e.g., "every," "some," "all") and the relationships between the different statements. Specifically, I will ask myself questions like: "Does the conclusion logically follow from the premises?" "Are there any implicit assumptions or unstated premises that could affect the validity of the argument?" "Are there any logical fallacies, such as affirming the consequent or denying the antecedent?" By asking these questions and carefully examining the arguments, I hope to improve my ability to distinguish deductively valid arguments from formal fallacies.} \\
\\ 
\texttt{To apply this strategy more effectively, I will also pay closer attention to the structure of the arguments, including the use of conditionals, disjunctions, and negations. I will also try to identify any potential ambiguities or unclear statements in the premises or conclusion, and ask myself how these might affect the validity of the argument.} \\
\\ 
\texttt{In particular, I will focus on the third argument, which I previously failed to identify as invalid. Upon re{-}examining the argument, I realize that I made a mistake in my previous analysis. The correct analysis is as follows:} \\
\\ 
\texttt{The premises state that every infrequent user of Paul Mitchell shampoo is either a rare consumer of Nioxin shampoo or a loyal buyer of Caress soap, or both. However, the conclusion states that whoever is an infrequent user of Paul Mitchell shampoo is not a regular consumer of Lush soap. This conclusion does not logically follow from the premises, as there is no information that infrequent users of Paul Mitchell shampoo are not regular consumers of Lush soap. In fact, it is possible that some infrequent users of Paul Mitchell shampoo are also regular consumers of Lush soap. Therefore, the argument is invalid.} \\
\\ 
\texttt{By carefully analyzing the premises and conclusion of this argument, I hope to improve my ability to identify formal fallacies and distinguish deductively valid arguments from invalid ones.} \\
\\ 
\texttt{To further improve my strategy, I will also try to identify the underlying logical principles that govern each argument, such as modus ponens, modus tollens, or hypothetical syllogism. By understanding these principles and how they are applied in each argument, I hope to develop a more nuanced and effective approach to evaluating the validity of arguments.} \\
\\ 
\\ 
\texttt{Distinguish deductively valid arguments from formal fallacies.} \\
\\ 
\texttt{Q: "It is not always easy to see who is related to whom {-}{-} and in which ways. The following argument pertains to this question: To begin with, Lesley is a close friend of Fernando. Moreover, being a close friend of Fernando or a schoolmate of Lowell is sufficient for being a great{-}grandfather of Leroy. It follows that Lesley is a great{-}grandfather of Leroy."} \\
\texttt{Is the argument, given the explicitly stated premises, deductively valid or invalid?} \\
\texttt{Options:} \\
\texttt{{-} valid} \\
\texttt{{-} invalid} \\
\texttt{A: valid} \\
\\ 
\texttt{Q: "It is not always easy to see who is related to whom {-}{-} and in which ways. The following argument pertains to this question: Whoever is not a great{-}grandfather of Clyde is a stepbrother of Brian. Being an ancestor of Dana is sufficient for not being a great{-}grandfather of Clyde. We may conclude: Everyone who is an ancestor of Dana is a stepbrother of Brian, too."} \\
\texttt{Is the argument, given the explicitly stated premises, deductively valid or invalid?} \\
\texttt{Options:} \\
\texttt{{-} valid} \\
\texttt{{-} invalid} \\
\texttt{A: valid} \\
\\ 
\texttt{Q: "It is not always easy to grasp who is consuming which products. The following argument pertains to this question: Every infrequent user of Paul Mitchell shampoo is either a rare consumer of Nioxin shampoo or a loyal buyer of Caress soap, or both. No regular consumer of Lush soap is a rare consumer of Nioxin shampoo and, in the same time, a loyal buyer of Caress soap. It follows that whoever is an infrequent user of Paul Mitchell shampoo is not a regular consumer of Lush soap."} \\
\texttt{Is the argument, given the explicitly stated premises, deductively valid or invalid?} \\
\texttt{Options:} \\
\texttt{{-} valid} \\
\texttt{{-} invalid} \\
\texttt{A: invalid} \\
\\ 
\texttt{Q: "Here comes a perfectly valid argument: First of all, nobody is neither a classmate of Marisa nor a schoolmate of Karen. Next, not every sister of Ruth is a schoolmate of Karen. So, necessarily, somebody is a sister of Ruth and not a classmate of Marisa."} \\
\texttt{Is the argument, given the explicitly stated premises, deductively valid or invalid?} \\
\texttt{Options:} \\
\texttt{{-} valid } \\
\texttt{{-} invalid} \\
\texttt{A:} \\
            }
        }
    }
     \\ \\

    \theutterance \stepcounter{utterance}  
    & \multicolumn{4}{p{0.6\linewidth}}{
        \cellcolor[rgb]{0.561,0.737,0.561}{
            \makecell[{{p{\linewidth}}}]{
                \texttt{\tiny{[P1$\rangle$GM]}}
                \texttt{Based on the provided plans and the new set of questions, I will apply the strategy to evaluate the arguments.} \\
\\ 
\texttt{1. "It is not always easy to see who is related to whom {-}{-} and in which ways. The following argument pertains to this question: To begin with, Lesley is a close friend of Fernando. Moreover, being a close friend of Fernando or a schoolmate of Lowell is sufficient for being a great{-}grandfather of Leroy. It follows that Lesley is a great{-}grandfather of Leroy."} \\
\\ 
\texttt{This argument is deductively valid. The premises state that Lesley is a close friend of Fernando and that being a close friend of Fernando or a schoolmate of Lowell is sufficient for being a great{-}grandfather of Leroy. Since Lesley is a close friend of Fernando, it logically follows that Lesley is a great{-}grandfather of Leroy.} \\
\\ 
\texttt{2. "It is not always easy to see who is related to whom {-}{-} and in which ways. The following argument pertains to this question: Whoever is not a great{-}grandfather of Clyde is a stepbrother of Brian. Being an ancestor of Dana is sufficient for not being a great{-}grandfather of Clyde. We may conclude: Everyone who is an ancestor of Dana is a stepbrother of Brian, too."} \\
\\ 
\texttt{This argument is deductively valid. The premises state that whoever is not a great{-}grandfather of Clyde is a stepbrother of Brian and that being an ancestor of Dana is sufficient for not being a great{-}grandfather of Clyde. Since being an ancestor of Dana implies not being a great{-}grandfather of Clyde, it logically follows that everyone who is an ancestor of Dana is a stepbrother of Brian.} \\
\\ 
\texttt{3. "It is not always easy to grasp who is consuming which products. The following argument pertains to this question: Every infrequent user of Paul Mitchell shampoo is either a rare consumer of Nioxin shampoo or a loyal buyer of Caress soap, or both. No regular consumer of Lush soap is a rare consumer of Nioxin shampoo and, in the same time, a loyal buyer of Caress soap. It follows that whoever is an infrequent user of Paul Mitchell shampoo is not a regular consumer of Lush soap."} \\
\\ 
\texttt{This argument is deductively invalid. The premises state that every infrequent user of Paul Mitchell shampoo is either a rare consumer of Nioxin shampoo or a loyal buyer of Caress soap, or both, and that no regular consumer of Lush soap is a rare consumer of Nioxin shampoo and a loyal buyer of Caress soap. However, the conclusion states that whoever is an infrequent user of Paul Mitchell shampoo is not a regular consumer of Lush soap. This conclusion does not logically follow from the premises, as there is no information that infrequent users of Paul Mitchell shampoo are not regular consumers of Lush soap. In fact, it is possible that some infrequent users of Paul Mitchell shampoo are also regular consumers of Lush soap.} \\
            }
        }
    }
    & & \\ \\

    \theutterance \stepcounter{utterance}  
    & & & \multicolumn{2}{p{0.3\linewidth}}{
        \cellcolor[rgb]{0.9,0.9,0.9}{
            \makecell[{{p{\linewidth}}}]{
                \texttt{\tiny{[GM$|$GM]}}
                \texttt{Target: invalid} \\
            }
        }
    }
    & & \\ \\

    \theutterance \stepcounter{utterance}  
    & & & \multicolumn{2}{p{0.3\linewidth}}{
        \cellcolor[rgb]{0.9,0.9,0.9}{
            \makecell[{{p{\linewidth}}}]{
                \texttt{\tiny{[GM$|$GM]}}
                \texttt{game\_result = LOSE} \\
            }
        }
    }
    & & \\ \\

    \theutterance \stepcounter{utterance}  
    & & \multicolumn{4}{p{0.6\linewidth}}{
        \cellcolor[rgb]{0.95,0.95,0.95}{
            \makecell[{{p{\linewidth}}}]{
                \texttt{\tiny{[P1$\leftrightarrow$P]}}
                \texttt{Post{-}episode Reflexion (not sent to GM):} \\
\texttt{New plan: } \\
\\ 
\texttt{Upon reflection, I realize that I failed to identify the formal fallacies in the arguments in my previous attempts. In particular, I did not recognize the invalidity of the third argument. To improve my strategy, I will focus on carefully analyzing the premises and the conclusion of each argument, looking for any logical flaws or inconsistencies. I will also pay closer attention to the use of quantifiers (e.g., "every," "some," "all") and the relationships between the different statements. Specifically, I will ask myself questions like: "Does the conclusion logically follow from the premises?" "Are there any implicit assumptions or unstated premises that could affect the validity of the argument?" "Are there any logical fallacies, such as affirming the consequent or denying the antecedent?" By asking these questions and carefully examining the arguments, I hope to improve my ability to distinguish deductively valid arguments from formal fallacies.} \\
\\ 
\texttt{To apply this strategy more effectively, I will also pay closer attention to the structure of the arguments, including the use of conditionals, disjunctions, and negations. I will also try to identify any potential ambiguities or unclear statements in the premises or conclusion, and ask myself how these might affect the validity of the argument.} \\
\\ 
\texttt{In particular, I will focus on the third argument, which I previously failed to identify as invalid. Upon re{-}examining the argument, I realize that I made a mistake in my previous analysis. The correct analysis is as follows:} \\
\\ 
\texttt{The premises state that every infrequent user of Paul Mitchell shampoo is either a rare consumer of Nioxin shampoo or a loyal buyer of Caress soap, or both. However, the conclusion states that whoever is an infrequent user of Paul Mitchell shampoo is not a regular consumer of Lush soap. This conclusion does not logically follow from the premises, as there is no information that infrequent users of Paul Mitchell shampoo are not regular consumers of Lush soap. In fact, it is possible that some infrequent users of Paul Mitchell shampoo are also regular consumers of Lush soap. Therefore, the argument is invalid.} \\
\\ 
\texttt{By carefully analyzing the premises and conclusion of this argument, I hope to improve my ability to identify formal fallacies and distinguish deductively valid arguments from invalid ones.} \\
\\ 
\texttt{To further improve my strategy, I will also try to identify the underlying logical principles that govern each argument, such as modus ponens, modus tollens, or hypothetical syllogism. By understanding these principles and how they are applied in each argument, I hope to develop a more nuanced and effective approach to evaluating the validity of arguments.} \\
\\ 
\texttt{Additionally, I will also try to identify any potential biases or assumptions that may be influencing my analysis. I will ask myself questions like: "Am I assuming something that is not explicitly stated in the premises?" "Am I relying on a particular interpretation of the language that may not be universally accepted?" By being more aware of these potential biases and assumptions, I hope} \\
            }
        }
    }
    & \\ \\

\end{supertabular}
}

\end{document}
